\subsection{Метод Зойтендейка}
\label{subsec:zoutendijk}

\[
  \varphi_0(x) \longrightarrow \min_{x \in S} ,
  \qquad
  S = \big\{\, x \in \R^n \ \big|\ \varphi_i(x) \le 0,\ i = \overline{1,m},
      \quad Ax = b \,\big\} ,
\]
где $A = A[l, n]$, $b = b[l]$, $l < n$.

\medskip
\noindent\textbf{Основная идея:} в точке $x \in S$ искать возможное направление
спуска (определение~\ref{def:feasible-direction}).

\subsubsection*{Предположения}

\begin{enumerate}[label=\arabic*),left=0pt]
  \item $\varphi_i(x)$ — выпуклые, $i = \overline{0,m}$;
  \item $\varphi_i(x)$ — непрерывно дифференцируемые, $i = \overline{0,m}$;
  \item $\norm{\nabla \varphi_i(x)} \le K$ \ $(\exists K > 0)$,
        $i = \overline{0,m}$;
  \item $\norm{\nabla\varphi_i(x) - \nabla\varphi_i(y)} \le R\norm{x - y}$,
        $R > 0$, $\forall x, y \in \R^n$;
  \item $S$ — компакт;
  \item $\exists\, \bar{x} \in S$: $\varphi_i(\bar{x}) < 0$, $i = \overline{1,m}$,
        $A\bar{x} = b$ — \textbf{условие Слейтера}.
\end{enumerate}

\subsubsection*{Активные и почти активные ограничения}

Пусть $M = \{1, 2, \dots, m\}$ и для $x \in S$
\[
  I_0(x) = \{\, i \in M \mid \varphi_i(x) = 0 \,\} ,
  \qquad
  I_1(x) = \{\, i \in M \mid \varphi_i(x) < 0 \,\} .
\]
Для $\delta > 0$ введём
\[
  I_\delta(x) = \{\, i \in M \mid -\delta \le \varphi_i(x) \le 0 \,\}
  \qquad \text{— почти активные ограничения.}
\]

\begin{idea}
Зачем нужен «зазор» $\delta$. Если учитывать только активные ограничения
$I_0(x)$, направление может оказаться возможным лишь формально: ограничение,
до которого остаётся ничтожно мало, тут же станет активным, и допустимая длина
шага окажется почти нулевой — метод заклинит у границы. Включая в рассмотрение
все ограничения, до которых меньше $\delta$, мы гарантируем шаг конечной длины.
Цена — направление получается более «осторожным»; поэтому $\delta$ по ходу
работы уменьшают (шаг~2$^\circ$ алгоритма ниже).
\end{idea}

\subsubsection*{Вспомогательная задача поиска направления}

Пусть $x \in S$, $\delta > 0$, $\xi_i > 0$, $i = \overline{0,m}$.
Вспомогательная задача:
\[
  \eta \longrightarrow \min_{\mathbb{H}} ,
\]
\[
  \mathbb{H} = \Big\{\, \eta \ \Big|\
    \nabla\T\varphi_i(x)\, s \le \xi_i\, \eta,\ i \in I_\delta(x) \cup \{0\} ;
    \quad A s = 0 ; \quad \norm{s}_\infty \le 1 \,\Big\} ,
\]
где $\norm{s}_\infty \le 1$ означает
$-1 \le s^{(i)} \le 1$, $i = \overline{1,n}$.

При $s = 0$, $\eta = 0$ все условия выполнены — пара допустима, поэтому решение
$s_\delta$, $\eta_\delta$ вспомогательной задачи удовлетворяет
\[
  \eta_\delta \le 0 .
\]

\begin{idea}
Переменная $\eta$ — это «насколько уверенно» направление $s$ одновременно
уменьшает цель и уводит от всех почти активных ограничений: неравенства
$\nabla\T\varphi_i(x)s \le \xi_i\eta$ при $\eta < 0$ означают, что по направлению
$s$ убывает и $\varphi_0$, и каждое почти активное $\varphi_i$. Минимизация
$\eta$ — поиск наилучшего такого направления.

Условие $As = 0$ держит точку на плоскости $Ax = b$, а $\norm{s}_\infty \le 1$
нормирует задачу: без него при $\eta < 0$ можно было бы бесконечно растягивать
$s$ и $\eta \to -\infty$. Коэффициенты $\xi_i > 0$ — веса, позволяющие
по-разному взвешивать цель и ограничения. Вся задача линейна по $(\eta, s)$,
то есть решается симплекс-методом из~\S1.

Случай $\eta_\delta = 0$ означает, что возможного направления спуска нет — это
и будет признаком оптимальности.
\end{idea}

\subsubsection*{Лемма о длине шага}

\begin{lemma}\label{lem:zoutendijk-step}
Если $\eta_\delta(x) < 0$, то существует $\alpha > 0$ такое, что
\begin{equation}\label{eq:zout-goal}
  \begin{cases}
    \varphi_0\big(x + \alpha s_\delta(x)\big)
      \le \varphi_0(x) + \tfrac12 \alpha\, \xi_0\, \eta_\delta(x) , \\[4pt]
    \varphi_i\big(x + \alpha s_\delta(x)\big) \le 0 , & i \in I_\delta(x) , \\[4pt]
    \varphi_i\big(x + \alpha s_\delta(x)\big) \le 0 , & i \notin I_\delta(x) .
  \end{cases}
\end{equation}
То есть $s_\delta(x)$ — возможное направление спуска.
\end{lemma}

\begin{proof}
По формуле конечных приращений
\[
  \varphi_0\big(x + \alpha s_\delta(x)\big)
  = \varphi_0(x) + \alpha\,\nabla\T\varphi_0(\tilde{x})\, s_\delta(x) ,
  \qquad
  \tilde{x} = x + \theta\,\alpha\, s_\delta(x) , \quad \theta \in [0,1] .
\]
Прибавим и вычтем $\alpha \nabla\T\varphi_0(x) s_\delta(x)$:
\[
  \varphi_0\big(x + \alpha s_\delta(x)\big)
  = \varphi_0(x) + \alpha\,\nabla\T\varphi_0(x)\, s_\delta(x)
    + \alpha\big(\nabla\varphi_0(\tilde{x}) - \nabla\varphi_0(x)\big)\T s_\delta(x) .
\]
Последнее слагаемое оцениваем неравенством Коши--Буняковского и условием~4):
\[
  \alpha\,\norm{\nabla\varphi_0(\tilde{x}) - \nabla\varphi_0(x)}\cdot
  \norm{s_\delta(x)}
  \le \alpha R \underbrace{\norm{\tilde{x} - x}}_{=\ \theta\alpha\norm{s_\delta(x)}}
      \norm{s_\delta(x)}
  \le \alpha^2 R\, \norm{s_\delta(x)}^2 .
\]
Наконец, по построению множества $\mathbb{H}$ выполнено
$\nabla\T\varphi_0(x)\, s_\delta(x) \le \xi_0\, \eta_\delta(x)$, откуда
\begin{equation}\label{eq:zout-2}
  \varphi_0\big(x + \alpha s_\delta(x)\big)
  \ \le\ \varphi_0(x) + \alpha\,\xi_0\,\eta_\delta(x)
         + \alpha^2 R\,\norm{s_\delta(x)}^2 .
\end{equation}

Для $i \in I_\delta(x)$ точно та же выкладка даёт
\begin{equation}\label{eq:zout-3}
  \varphi_i\big(x + \alpha s_\delta(x)\big)
  \ \le\ \varphi_i(x) + \alpha\,\xi_i\,\eta_\delta(x)
         + \alpha^2 R\,\norm{s_\delta(x)}^2 .
\end{equation}

Для $i \notin I_\delta(x)$, то есть $\varphi_i(x) < -\delta$, используем
условие~3):
\begin{equation}\label{eq:zout-4}
  \begin{aligned}
    \varphi_i\big(x + \alpha s_\delta(x)\big)
    &= \varphi_i(x) + \alpha\,\nabla\T\varphi_i(\tilde{x})\, s_\delta(x) \\
    &\le \varphi_i(x)
       + \alpha\,\norm{\nabla\varphi_i(\tilde{x})}\cdot\norm{s_\delta(x)}
     \ <\ -\delta + \alpha K \norm{s_\delta(x)} .
  \end{aligned}
\end{equation}

Теперь подберём $\alpha > 0$ так, чтобы выполнялось~\eqref{eq:zout-goal}.
\begin{itemize}[left=0pt]
  \item Из~\eqref{eq:zout-2} достаточно
        $\alpha\xi_0\eta_\delta(x) + \alpha^2 R\norm{s_\delta(x)}^2
         \le \tfrac12 \alpha \xi_0 \eta_\delta(x)$, то есть
        \begin{equation}\label{eq:zout-6}
          \alpha \ \le\ -\frac{1}{2}\cdot
          \frac{\xi_0\,\eta_\delta(x)}{R\,\norm{s_\delta(x)}^2} .
        \end{equation}
  \item Из~\eqref{eq:zout-3} при $i \in I_\delta(x)$ имеем $\varphi_i(x) \le 0$,
        поэтому достаточно
        $\alpha\xi_i\eta_\delta(x) + \alpha^2 R\norm{s_\delta(x)}^2 \le 0$, то есть
        \begin{equation}\label{eq:zout-7}
          \alpha \ \le\ -\frac{\xi_i\,\eta_\delta(x)}{R\,\norm{s_\delta(x)}^2} .
        \end{equation}
  \item Из~\eqref{eq:zout-4} при $i \notin I_\delta(x)$ достаточно
        $-\delta + \alpha K \norm{s_\delta(x)} \le 0$, то есть
        \begin{equation}\label{eq:zout-8}
          \alpha \ \le\ \frac{\delta}{K\,\norm{s_\delta(x)}} .
        \end{equation}
\end{itemize}
Так как $\eta_\delta(x) < 0$, правые части
\eqref{eq:zout-6}--\eqref{eq:zout-8} положительны, и годится любое
$\alpha > 0$, не превосходящее их минимума.

Равенства сохраняются при любом $\alpha$:
\[
  A\big(x + \alpha s_\delta(x)\big)
  = Ax + \alpha \underbrace{A s_\delta(x)}_{=\,0} = b .
  \qedhere
\]
\end{proof}

\begin{idea}
Три оценки имеют разный смысл. Условие~\eqref{eq:zout-6} требует, чтобы шаг был
достаточно коротким и квадратичная поправка $\alpha^2 R\norm{s_\delta}^2$ не
съела линейное убывание — остаётся хотя бы его половина. Условие
\eqref{eq:zout-7} не даёт выйти за почти активные ограничения, а
\eqref{eq:zout-8} — за «далёкие»: до них запас ровно $\delta$, а скорость
изменения $\varphi_i$ не больше $K\norm{s_\delta}$.

Именно здесь работает весь набор предположений: липшицевость градиентов~4) даёт
квадратичную поправку, ограниченность градиентов~3) — оценку
\eqref{eq:zout-8}, а $\eta_\delta < 0$ — положительность всех трёх границ.
\end{idea}

\begin{theorem}\label{thm:zoutendijk}
Если $x \ne x_*$, $x \in S$, то $\eta_\delta(x) < 0$ для достаточно малых
$\delta > 0$.
\end{theorem}

\begin{idea}
Содержательно: пока точка не оптимальна, возможное направление спуска
существует — надо лишь взять достаточно узкую полосу почти активных
ограничений. Вместе с леммой~\ref{lem:zoutendijk-step} это и даёт сходимость:
на каждом шаге либо есть направление с $\eta_k < 0$ и целевая функция
убывает на конечную величину, либо $\delta$ уменьшается. Условие Слейтера~6)
нужно именно здесь — без внутренней точки конус возможных направлений может
выродиться.
\end{idea}

\subsubsection*{Алгоритм}

Пусть заданы $\xi_0$ и $\{\xi_i\}_{i=1}^{m}$: $\xi_i > 0$ для всех
$i = \overline{0,m}$; $x_0$ — начальное приближение, $\delta_0 > 0$,
$\lambda \in (0,1)$.

\medskip
\noindent\textbf{Общий шаг в точке $x_k$.}

\begin{steps}
  \item \textbf{Поиск возможного направления спуска:}
        $\eta \to \min_{\mathbb{H}}$, где
        \[
          \mathbb{H} = \Big\{\, \eta \ \Big|\
            \nabla\T\varphi_0(x_k)\, s \le \xi_0 \eta ; \quad
            \nabla\T\varphi_i(x_k)\, s \le \xi_i \eta,\ i \in I_{\delta_k}(x_k) ;
        \]
        \[
            A s = 0 ; \quad \norm{s}_\infty \le 1 \,\Big\} .
        \]
        Решение обозначаем $\eta_{\delta_k} \equiv \eta_k$,
        $s_{\delta_k} \equiv s_k$.
  \item \textbf{Выбор шага, построение $x_{k+1}$, $\delta_{k+1}$:}
        \begin{itemize}[left=0pt,nosep]
          \item[$1^\circ$] $\eta_k < -\delta_k$ — далеко от границы, двигаемся:
                \[
                  x_{k+1} = x_k + \alpha_k s_k , \qquad \delta_{k+1} = \delta_k ,
                \]
                где $\alpha_k$ выбирается по принципу дробления;
          \item[$2^\circ$] $0 > \eta_k \ge -\delta_k$ — направление слишком
                слабое, сужаем полосу:
                \[
                  x_{k+1} = x_k , \qquad \delta_{k+1} = \lambda\,\delta_k .
                \]
        \end{itemize}
        В обоих случаях $k := k+1$ и возврат к шагу~1.
\end{steps}

\medskip
\noindent\textbf{Условие останова:}
\[
  \delta_k < \delta_{0k} ,
  \qquad
  \delta_{0k} = -\max_{i \notin I_0(x_k)} \varphi_i(x_k) ,
\]
и $\eta_k = 0$. Тогда $x_k$ — решение.

\begin{idea}
Величина $\delta_{0k}$ — расстояние (по значению функций) до ближайшего
\emph{неактивного} ограничения: для $i \notin I_0(x_k)$ все $\varphi_i(x_k) < 0$,
и $-\max$ по ним даёт наименьший запас. Неравенство $\delta_k < \delta_{0k}$
означает, что полоса почти активных ограничений сузилась настолько, что
$I_{\delta_k}(x_k) = I_0(x_k)$ — в рассмотрении остались только по-настоящему
активные ограничения. Если при этом $\eta_k = 0$, то возможного направления
спуска нет даже в точном смысле, и $x_k$ оптимальна.

Логика всего алгоритма: пункт $1^\circ$ продвигает точку, пункт $2^\circ$
уточняет модель. Каждый проход через $2^\circ$ уменьшает $\delta$ в $\lambda$
раз, так что либо точка движется с ощутимым убыванием $\varphi_0$, либо
$\delta_k \to 0$ и срабатывает условие останова.
\end{idea}

\subsubsection*{Как найти начальное приближение}

Решается вспомогательная задача
\[
  \eta \longrightarrow \min_{\eta \in \mathbb{H}} ,
  \qquad
  \mathbb{H} = \big\{\, \eta \ \big|\ \varphi_i(x) - \eta \le 0,\
    i = \overline{1,m} \ \ \&\ \ Ax = b \,\big\} .
\]
По условию Слейтера существует $\bar{x}$: $\varphi_i(\bar{x}) < 0$ и
$A\bar{x} = b$, поэтому $\min \eta < 0$. Берём приближение $x$, удовлетворяющее
$Ax = b$.

\begin{idea}
Это та же идея, что и у вспомогательной задачи симплекс-метода (\S4): вместо
поиска допустимой точки «наугад» минимизируется максимальная невязка
$\eta = \max_i \varphi_i(x)$. Отрицательное значение минимума означает, что
найденная точка строго внутри всех неравенств — ровно то, что нужно для старта
метода, работающего с почти активными ограничениями.
\end{idea}
